% GENERATED FILE. Edit canonical lessons, then run npm run generate:books.
% canonical-source-hash: fnv1a64:36d45251ca8e0f6f
% canonical-lessons: UR-C176-naukri, UR-C176-pesha, UR-C176-tankhvah, UR-C176-rupaya, UR-C176-sikka

\chapter{Job, Profession, Salary, Rupee, Coin}
\label{ch:ur-job-profession-salary-rupee-coin}
\hlchaptermodality{176}

\begin{chapteropening}
\textbf{By the end of this chapter:} \emph{I can say a job, a profession, a salary, a rupee and a coin.}

Complete the last lesson of chapter 176: I can say a job, a profession, a salary, a rupee and a coin.
\end{chapteropening}

\section[\texorpdfstring{\ur{نوکری} (naukrī)}{نوکری (naukrī)}]{\ur{نوکری} (naukrī) — a job}

\label{lesson:UR-C176-naukri}



\begin{quote}
Before the new one: say the Urdu for a quail, then the Urdu for a dove.
\end{quote}

\subsection*{You'll want to know: \ur{نوکری}}
\textbf{\ur{نوکری}} — \emph{naukrī} — \textquotedblleft{}a job\textquotedblright{}.

\begin{cousinweb}
\textbf{\ur{نوکری}} (\emph{naukrī}) is \emph{naukar}, \textquotedblleft{}servant\textquotedblright{}, with the ending \emph{-ī} that makes a noun. \emph{naukar} is usually traced to Persian \emph{naukar}, which Persian took from Mongolic \emph{nökör}, \textquotedblleft{}companion\textquotedblright{}.
\end{cousinweb}

\begin{grammarlens}[title={What you've built}]
One more word for this chapter.
\end{grammarlens}

\subsection*{Your turn}
\begin{itemize}
\raggedright
  \item \emph{Say it:} \emph{naukrī}
  \item \emph{Say it:} \emph{naukrī}, once more
  \item \emph{Say it:} say \emph{naukrī}
  \item \emph{Recall:} say the Urdu for a quail, then the Urdu for a dove, then say \emph{naukrī} again
\end{itemize}

\begin{tcolorbox}[breakable,colback=teal!4,colframe=teal!35!black,title={Before you move on}]
What does \ur{نوکری} mean? (\textquotedblleft{}A job\textquotedblright{}.) Say it once more.
\end{tcolorbox}
\section[\texorpdfstring{\ur{پیشہ} (peshā)}{پیشہ (peshā)}]{\ur{پیشہ} (peshā) — a profession}

\label{lesson:UR-C176-pesha}



\begin{quote}
Before the new one: say the Urdu for a dove, then the Urdu for a job.
\end{quote}

\subsection*{You'll want to know: \ur{پیشہ}}
\textbf{\ur{پیشہ}} — \emph{peshā} — \textquotedblleft{}a profession\textquotedblright{}.

\begin{cousinweb}
\textbf{\ur{پیشہ}} (\emph{peshā}) is a loan, usually traced to Persian \emph{pesha}, \textquotedblleft{}craft, profession\textquotedblright{}.
\end{cousinweb}

\begin{grammarlens}[title={What you've built}]
One more word for this chapter.
\end{grammarlens}

\subsection*{Your turn}
\begin{itemize}
\raggedright
  \item \emph{Say it:} \emph{peshā}
  \item \emph{Say it:} \emph{peshā}, once more
  \item \emph{Say it:} say \emph{peshā}
  \item \emph{Recall:} say the Urdu for a dove, then the Urdu for a job, then say \emph{peshā} again
\end{itemize}

\begin{tcolorbox}[breakable,colback=teal!4,colframe=teal!35!black,title={Before you move on}]
What does \ur{پیشہ} mean? (\textquotedblleft{}A profession\textquotedblright{}.) Say it once more.
\end{tcolorbox}
\section[\texorpdfstring{\ur{تنخواہ} (tankhvāh)}{تنخواہ (tankhvāh)}]{\ur{تنخواہ} (tankhvāh) — a salary}

\label{lesson:UR-C176-tankhvah}



\begin{quote}
Before the new one: say the Urdu for a job, then the Urdu for a profession.
\end{quote}

\subsection*{You'll want to know: \ur{تنخواہ}}
\textbf{\ur{تنخواہ}} — \emph{tankhvāh} — \textquotedblleft{}a salary\textquotedblright{}.

\begin{cousinweb}
\textbf{\ur{تنخواہ}} (\emph{tankhvāh}) is a loan, usually traced to Persian \emph{tankhwāh}, \textquotedblleft{}salary\textquotedblright{}.
\end{cousinweb}

\begin{grammarlens}[title={What you've built}]
One more word for this chapter.
\end{grammarlens}

\subsection*{Your turn}
\begin{itemize}
\raggedright
  \item \emph{Say it:} \emph{tankhvāh}
  \item \emph{Say it:} \emph{tankhvāh}, once more
  \item \emph{Say it:} say \emph{tankhvāh}
  \item \emph{Recall:} say the Urdu for a job, then the Urdu for a profession, then say \emph{tankhvāh} again
\end{itemize}

\begin{tcolorbox}[breakable,colback=teal!4,colframe=teal!35!black,title={Before you move on}]
What does \ur{تنخواہ} mean? (\textquotedblleft{}A salary\textquotedblright{}.) Say it once more.
\end{tcolorbox}
\section[\texorpdfstring{\ur{روپیہ} (rupayā)}{روپیہ (rupayā)}]{\ur{روپیہ} (rupayā) — a rupee}

\label{lesson:UR-C176-rupaya}



\begin{quote}
Before the new one: say the Urdu for a profession, then the Urdu for a salary.
\end{quote}

\subsection*{You'll want to know: \ur{روپیہ}}
\textbf{\ur{روپیہ}} — \emph{rupayā} — \textquotedblleft{}a rupee\textquotedblright{}.

\begin{cousinweb}
\textbf{\ur{روپیہ}} (\emph{rupayā}) is usually traced to Sanskrit \emph{rūpya}, \textquotedblleft{}bearing a stamp; silver\textquotedblright{} (CDIAL 10805). CDIAL lists the same word in Hindi as \emph{rūpā}, \textquotedblleft{}silver\textquotedblright{}.
\end{cousinweb}

\begin{grammarlens}[title={What you've built}]
One more word for this chapter.
\end{grammarlens}

\subsection*{Your turn}
\begin{itemize}
\raggedright
  \item \emph{Say it:} \emph{rupayā}
  \item \emph{Say it:} \emph{rupayā}, once more
  \item \emph{Say it:} say \emph{rupayā}
  \item \emph{Recall:} say the Urdu for a profession, then the Urdu for a salary, then say \emph{rupayā} again
\end{itemize}

\begin{tcolorbox}[breakable,colback=teal!4,colframe=teal!35!black,title={Before you move on}]
What does \ur{روپیہ} mean? (\textquotedblleft{}A rupee\textquotedblright{}.) Say it once more.
\end{tcolorbox}
\section[\texorpdfstring{\ur{سکہ} (sikkā)}{سکہ (sikkā)}]{\ur{سکہ} (sikkā) — a coin}

\label{lesson:UR-C176-sikka}



\begin{quote}
Before the new one: say the Urdu for a salary, then the Urdu for a rupee.
\end{quote}

\subsection*{You'll want to know: \ur{سکہ}}
\textbf{\ur{سکہ}} — \emph{sikkā} — \textquotedblleft{}a coin\textquotedblright{}.

\begin{cousinweb}
\textbf{\ur{سکہ}} (\emph{sikkā}) is a loan from Arabic \emph{sikka}, \textquotedblleft{}a die for stamping coins\textquotedblright{}, which reached Urdu through Persian. Its Arabic root is \emph{s-k-k} (Lane).
\end{cousinweb}

\begin{grammarlens}[title={What you've built}]
One more word for this chapter.
\end{grammarlens}

\subsection*{Your turn}
\begin{itemize}
\raggedright
  \item \emph{Say it:} \emph{sikkā}
  \item \emph{Say it:} \emph{sikkā}, once more
  \item \emph{Say it:} say \emph{sikkā}
  \item \emph{Recall:} say the Urdu for a salary, then the Urdu for a rupee, then say \emph{sikkā} again
\end{itemize}

\begin{tcolorbox}[breakable,colback=teal!4,colframe=teal!35!black,title={Before you move on}]
What does \ur{سکہ} mean? (\textquotedblleft{}A coin\textquotedblright{}.) Say it once more.
\end{tcolorbox}
